\chapter{Seguridad y respuesta a incidentes}

\section{Controles implementados}

\begin{tabularx}{\textwidth}{@{}L{0.25\textwidth}Y@{}}
\toprule
\textbf{Riesgo} & \textbf{Control}\\
\midrule
Robo de sesión & access token en memoria; refresh en cookie `HttpOnly` y `Secure`\\
Acceso indebido & permisos por rol y consultas filtradas en backend\\
Abuso de API & throttling anónimo/autenticado y paginación\\
Origen no autorizado & CORS, CSRF y WebSocket limitados a orígenes exactos\\
Archivos inseguros & validación de tipo/tamaño y Storage solo desde backend\\
Secretos & variables de entorno, ejemplos con marcadores y archivos reales ignorados\\
Transporte & HTTPS, HSTS y cookies seguras fuera de `DEBUG`\\
\bottomrule
\end{tabularx}

Ocultar un botón en React no es un control de seguridad. Toda decisión autoritativa reside en
permisos, servicios y repositorios del backend.

\section{Manejo de secretos}

Use el gestor de secretos de cada plataforma. Entregue valores únicamente a personas autorizadas y
por un canal separado del código. Mantenga un inventario con propietario, propósito, última rotación
y procedimiento de revocación.

Nunca publique `.env`, claves Supabase, secretos Django, tokens, contraseñas de correo, exportaciones,
capturas con PII ni copias de base de datos. Si un secreto entra en Git, rotarlo es obligatorio;
borrarlo del último commit no invalida su historial.

\section{Datos demo}

`seed\_demo` requiere confirmación explícita y, con `DEBUG=False`, autorización adicional para un
staging aislado. Su propósito es demostración o prueba. Las cuentas generadas son activas y pueden
incluir privilegios administrativos; no habilite la excepción en producción.

\section{Reporte responsable}

No abra un issue público para una vulnerabilidad. Use un Security Advisory privado de GitHub o el
canal privado designado por SassBlum. Incluya componente, versión/SHA, reproducción mínima sin datos
reales, impacto y mitigación propuesta.

\section{Respuesta a incidentes}

\begin{enumerate}
  \item preserve logs, hora, entorno y SHA sin destruir evidencia;
  \item contenga el acceso y determine alcance;
  \item revoque sesiones y rote secretos comprometidos;
  \item corrija y añada prueba de regresión;
  \item valide en staging y despliegue con rollback;
  \item notifique según el procedimiento acordado;
  \item documente causa, tiempos y acciones preventivas.
\end{enumerate}

\ManualNote{Solo la rama `main` desplegada recibe soporte hasta que exista una release estable. La
política de versiones debe actualizarse al publicar el primer tag contractual.}
